\documentclass[11pt]{article}

\usepackage[letterpaper,margin=1in]{geometry}
\usepackage[T1]{fontenc}
\usepackage{lmodern}
\usepackage{amsmath}
\usepackage{amssymb}
\usepackage{amsthm}
\usepackage{enumitem}
\usepackage{longtable}
\usepackage{booktabs}
\usepackage{natbib}
\usepackage[hidelinks]{hyperref}

\newtheorem{proposition}{Proposition}
\newtheorem{theorem}{Theorem}
\newtheorem{corollary}{Corollary}
\theoremstyle{definition}
\newtheorem{definition}{Definition}

\setlist[enumerate]{itemsep=1pt,topsep=3pt}

\newcommand{\cN}{\mathcal N}
\newcommand{\cK}{\mathcal K}
\newcommand{\cG}{\mathcal G}
\newcommand{\cP}{\mathcal P}
\newcommand{\cA}{\mathcal A}
\newcommand{\TR}{\mathit{TR}}

\begin{document}

\section{Model}
\label{sec:model}

We develop a static general equilibrium model of multinational production in which
the strategic agent is the ownership group. A group can operate plants in several
countries and, when it sets quantities, takes into account that its own plants compete
with one another. Markups therefore depend on the group's total share of the market
rather than on the share of any single plant. Profits are paid to the countries that
own each group, and ownership is recorded group by group rather than country by
country. The baseline takes as given which plants each group operates and which
markets they serve; Section~\ref{sec:entry} lets groups choose them. Proofs are in the
Appendix, and a table of notation closes the document.

\subsection{Model Set-Up}
\label{sec:setup}

\paragraph{Agents.} There is a finite set of countries $\cN=\{1,\dots,N\}$ and a
finite set of sectors $\cK$. Each country $n$ is inhabited by a representative
household that supplies $L_n$ units of labor inelastically and owns shares in firms.
Production takes place in plants, and every plant belongs to exactly one group
$g\in\cG$, which we identify with its global ultimate parent. Each group has a
headquarters country $h(g)$, operates in a single sector $k(g)$,\footnote{This
restriction is without loss for behavior. Marginal costs are constant, so a group's
decisions separate across markets (Section~\ref{sec:market}). A group active in
several sectors is equivalent to several single-sector groups with the same owners;
what the restriction rules out is a capability shared across sectors.} and owns a
finite set of plants $\cP_g$. We index plants by $i$ and write $g(i)$ and $\ell(i)$ for
the group and the location of plant $i$. Each plant produces one differentiated
variety.

\paragraph{Representative household.} The household in country $n$ has
Cobb--Douglas preferences across sectors and CES preferences within them,
\begin{equation}
  U_n=\prod_{k\in\cK}C_{nk}^{\beta_k},
  \qquad
  C_{nk}=\Bigl[\sum_{i\in\cA_{nk}}q_{i,n}^{(\sigma_k-1)/\sigma_k}\Bigr]^{\sigma_k/(\sigma_k-1)},
  \label{eq:utility}
\end{equation}
where $\cA_{nk}$ is the set of plants that sell sector-$k$ goods in country $n$,
$\beta_k>0$, $\sum_k\beta_k=1$, and $\sigma_k>1$ is the elasticity of substitution
between varieties of the same sector. The household spends the share $\beta_k$ of its
income $X_n$ on sector $k$. Writing $q_{i,n}$ for the quantity of plant $i$'s variety
sold in $n$ and $E_{nk}$ for spending on sector $k$ in $n$, demand is
\begin{equation}
  q_{i,n}=p_{i,n}^{-\sigma_k}\,P_{nk}^{\sigma_k-1}E_{nk},
  \qquad
  P_{nk}=\Bigl[\sum_{i\in\cA_{nk}}p_{i,n}^{1-\sigma_k}\Bigr]^{1/(1-\sigma_k)} ,
  \label{eq:demand}
\end{equation}
where $P_{nk}$ is the sector's price index. Welfare is real income $X_n/P_n$, with
$P_n=\prod_k(P_{nk}/\beta_k)^{\beta_k}$.\footnote{With Cobb--Douglas preferences across
sectors, a sector's spending does not expand when its price index falls. A CES upper
tier with elasticity $\eta\in(1,\sigma_k)$ would add that margin and change the markup
in \eqref{eq:markup} to $1/\mu(S)=1-(1-S)/\sigma_k-S/\eta$; the structure of the market
equilibrium below is unchanged.}

\paragraph{Production technology.} Labor is the only factor of production. A plant
located in $\ell$ produces with constant returns, and its unit cost is
\begin{equation}
  a_i=\frac{w_\ell}{\varphi_i},
  \label{eq:cost}
\end{equation}
where $w_\ell$ is the wage in $\ell$ and $\varphi_i>0$ the plant's productivity.%
\footnote{We do not include intermediate inputs. With an input--output structure, a
tariff also raises the cost of the input bundle, and the existence and uniqueness
results below require additional conditions.}

\paragraph{Head-office capability.} Productivity can combine a component of the
location, a component of the plant, and a component common to all plants of the
group. The group component captures head-office services such as management,
technology and brands \citep{Antras2003,AntrasHelpman2004}. These services are a
non-rival capability of the group that all of its plants share, so they enter
productivity rather than the cost of inputs.\footnote{The common alternative treats
head-office services as labor bought at the parent's wage, so that a plant abroad pays
part of its costs at a foreign wage. This links factor markets across countries and can
overturn the substitutability of labor across countries on which uniqueness of wages
rests.} A plant may also be less productive when it operates away from its group's
headquarters \citep{RamondoRodriguezClare2013,HeadMayer2019}. The parametric form of
$\varphi_i$ matters for the quantification, not for the results of this section.

\paragraph{Trade.} Shipping goods from $\ell$ to $n$ entails an iceberg cost
$d_{\ell n}\ge1$, with $d_{nn}=1$. Trade costs run from the country where the good is
produced, not from headquarters, so a single plant can serve several destinations,
including third markets, as an export platform \citep{Tintelnot2017}. Country $n$
levies a tariff $t_{\ell nk}\ge0$ on sector-$k$ goods produced in $\ell$, with
$t_{nnk}=0$. Tariffs depend on where a good is produced, not on who owns the plant that
produces it, so a country's tariffs also fall on the foreign affiliates of the groups
its residents own. We model the tariff as a proportional wedge on the cost of
delivering the good. The delivered marginal cost of plant $i$ in market $(n,k)$ is
\begin{equation}
  c_{i,n}=a_i\,d_{\ell n}\,(1+t_{\ell nk}),\qquad \ell=\ell(i),
  \label{eq:delivered}
\end{equation}
and country $n$ collects $t_{\ell nk}\,a_id_{\ell n}q_{i,n}$ on the sales of plant
$i$.\footnote{A duty on the exporter's price yields the same prices when all the plants
a group uses to serve $n$ face the same tariff. When they do not, it makes the group's
markup differ across its plants, and the property that a group's markup depends only on
its total share (Section~\ref{sec:market}) is lost.}

\paragraph{Multinational corporations.} A \emph{market} is a destination--sector pair
$(n,k)$. In the baseline, each plant serves a given set of destinations, so the set
$\cA_{nk}$ of plants that serve market $(n,k)$ is given. A group with plants in more
than one country is a multinational, a plant located outside its group's headquarters
country is a foreign affiliate, and a group with a single plant is a stand-alone firm.
A multinational can serve a market from a plant at home, from a plant in the
destination, from a plant in a third country, or from several at once. We assume that
every market is served by at least two groups. Section~\ref{sec:entry} lets groups
choose which plants serve each market, so that multinational status becomes an
equilibrium outcome.

\paragraph{Profit repatriation.} Country $n$ owns a share $\theta_{gn}\ge0$ of group
$g$, with $\sum_n\theta_{gn}=1$. Groups maximize their own profits whoever owns them:
ownership determines where profits are spent, not how groups price. The profit of a
group, $\pi_g=\sum_n\pi_{g,n}$, where $\pi_{g,n}$ is its profit in market $(n,k(g))$
(Section~\ref{sec:market}), is paid to its owners in proportion to $\theta_{gn}$,
wherever it is earned. Tariff revenue is rebated lump sum to the household of the
country that collects it. Income in country $n$ is
\begin{equation}
  X_n=w_nL_n+\sum_{g\in\cG}\theta_{gn}\,\pi_g+\TR_n,
  \qquad
  \TR_n=\sum_{k}\sum_{i\in\cA_{nk}}t_{\ell(i)nk}\,a_i\,d_{\ell(i)n}\,q_{i,n}.
  \label{eq:income}
\end{equation}
Profits are earned where groups sell and received where their owners reside, and
$\theta$ links the two. A natural benchmark assigns each group to the country of its
ultimate parent, $\theta_{gn}=\mathbf 1\{n=h(g)\}$; nothing below requires this form.
Because income includes the profits of the groups a country owns, wherever they
produce and sell, trade need not balance country by country.

\paragraph{Timing.} The model is static. Plants, productivities and ownership shares
are given. Groups choose the quantities their plants sell in each market, taking
wages and market sizes as given and taking into account their effect on the market's
price index, and households choose their consumption.

\subsection{Market Structure}
\label{sec:market}

Fix a market $(n,k)$ and suppress its indices: we write $\sigma$, $E$ and $P$ for
$\sigma_k$, $E_{nk}$ and $P_{nk}$, $c_i$ for the delivered cost $c_{i,n}$, and $T_g$ for
the plants of group $g$ that serve the market.

\paragraph{Demand.} Let $A\equiv\sum_{i\in\cA}p_i^{1-\sigma}=P^{1-\sigma}$. By
\eqref{eq:demand}, the revenue share of plant $i$ and the share of group $g$ are
\begin{equation}
  s_i=\frac{p_iq_i}{E}=\frac{p_i^{1-\sigma}}{A},
  \qquad
  S_g=\sum_{i\in T_g}s_i,
  \qquad
  \sum_g S_g=1 .
  \label{eq:shares}
\end{equation}
The aggregate $A$ summarizes how competitive the market is: it is large when rivals
charge low prices, and then any given price wins a small share.

\paragraph{Competition.} Groups are the strategic agents. Each group chooses the
quantities of all of its plants that serve the market, taking as given the quantities
of other groups and the level of spending $E$. It maximizes its profit
$\pi_g=\sum_{i\in T_g}(p_i-c_i)q_i$.

\begin{proposition}[Markups]\label{prop:markup}
The first-order conditions of group $g$ are $p_i=\mu(S_g)\,c_i$ for every $i\in T_g$,
with
\begin{equation}
  \mu(S)=\frac{\sigma}{(\sigma-1)(1-S)} .
  \label{eq:markup}
\end{equation}
\end{proposition}

\noindent This is the markup of \citet{AtkesonBurstein2008}, evaluated at the share
of the group rather than of the plant.\footnote{Under price competition the markup is
$\varepsilon(S)/[\varepsilon(S)-1]$ with $\varepsilon(S)=\sigma-(\sigma-1)S$. It is also
common to the group's plants and increasing in its share, so the structure below
carries over, but profits are no longer quadratic in the share and the ownership result
of Section~\ref{sec:cs} becomes an approximation.} It rises from $\sigma/(\sigma-1)$ for
a negligible group and is unbounded as the share tends to one, which is why every
market needs at least two groups.

\paragraph{Cannibalization across multinational affiliates.} When a group expands one
of its plants, the sector price index falls and every variety in the market loses
sales, the group's other plants included. A single-plant firm ignores the losses it
imposes on others; a group counts the losses it imposes on itself, and these grow with
its share. This is the factor $1-S_g$ in \eqref{eq:markup}. Two implications follow.
First, all plants of a group charge the same markup in a given market, and differences
in their costs show up in quantities rather than in markups. Second, market power is a
property of the group: concentration should be measured at the level of groups, and a
group charges more than its plants would charge as independent firms. Within a group,
\eqref{eq:shares} and the common markup give
\begin{equation}
  \frac{s_i}{S_g}=\frac{c_i^{1-\sigma}}{K_g},
  \qquad
  K_g\equiv\sum_{i\in T_g}c_i^{1-\sigma}.
  \label{eq:strength}
\end{equation}
We call $K_g$ the \emph{strength} of the group in the market: an index of its cost
efficiency that rises with the number of plants serving the market and with how cheaply
each delivers. Substituting $p_i=\mu(S_g)c_i$ into \eqref{eq:shares},
\begin{equation}
  \psi(S_g)=\frac{K_g}{A},
  \qquad
  \psi(S)\equiv S\,\mu(S)^{\sigma-1}
  =\Bigl(\frac{\sigma}{\sigma-1}\Bigr)^{\sigma-1}\frac{S}{(1-S)^{\sigma-1}} .
  \label{eq:sharemap}
\end{equation}
The map $\psi$ is strictly increasing, so $S_g=\psi^{-1}(K_g/A)$: the share of a group
depends only on its strength relative to the competitiveness of the market. However
many plants a group runs, and in however many countries, its position in the market is
summarized by one number, $K_g$, and all it needs to know about its rivals is $A$. The
equilibrium of the market is the value of $A$ at which shares add up to one,
\begin{equation}
  \sum_{g}\psi^{-1}\!\left(\frac{K_g}{A}\right)=1 .
  \label{eq:clearing}
\end{equation}

\paragraph{Profits.} Since $c_iq_i=p_iq_i/\mu(S_g)$, the profit of group $g$ is
$\pi_g=E\,S_g\,m(S_g)$, where $m(S)\equiv1-1/\mu(S)$ is the Lerner index. By
\eqref{eq:markup} the Lerner index is linear in the share, $m(S)=(1+\kappa S)/\sigma$
with $\kappa\equiv\sigma-1$, so that
\begin{equation}
  \pi_g=\frac{E}{\sigma}\bigl(S_g+\kappa S_g^2\bigr).
  \label{eq:profit}
\end{equation}
Profits rise more than proportionally with the share, because a larger group both sells
more and earns more on each unit sold. The curvature $\kappa$ governs the ownership
results of Section~\ref{sec:cs}. It would be zero with constant markups.

\subsection{General Equilibrium}
\label{sec:ge}

\paragraph{Market sizes.} Without intermediate inputs, all spending is final spending
by households, so spending on market $(n,k)$ is
\begin{equation}
  E_{nk}=\beta_kX_n .
  \label{eq:spending}
\end{equation}

\paragraph{Labor market.} Let $V_{i,n'}\equiv a_id_{\ell(i)n'}q_{i,n'}$ be the cost of
plant $i$ in serving country $n'$, exclusive of tariffs. By
Proposition~\ref{prop:markup}, $V_{i,n'}=r_{i,n'}/[\mu_{g(i)}(1+t_{\ell(i)n'k})]$, where
$r_{i,n'}$ is the plant's revenue in that market, $\mu_{g(i)}$ its group's markup
there, and $k$ the plant's sector. All of this cost is paid to labor where the plant is
located, so labor markets clear when
\begin{equation}
  w_nL_n=\sum_{k\in\cK}\sum_{n'\in\cN}\sum_{i\in\cA_{n'k}:\,\ell(i)=n}V_{i,n'} .
  \label{eq:labor}
\end{equation}

\begin{definition}[Equilibrium]\label{def:eq}
An equilibrium consists of wages $\{w_n\}$, incomes $\{X_n\}$, and prices and
quantities $\{p_{i,n},q_{i,n}\}$ such that households maximize \eqref{eq:utility}; in
every market, prices satisfy Proposition~\ref{prop:markup} and quantities satisfy
\eqref{eq:demand}; incomes satisfy \eqref{eq:income}, with market sizes
\eqref{eq:spending}; and labor markets clear, \eqref{eq:labor}.
\end{definition}

\noindent Prices, costs and incomes are homogeneous of degree one in wages, and shares
are homogeneous of degree zero, so wages are determined up to scale. At any wages, the
value of labor demanded across all countries equals the value of labor supplied
(Walras' law, Appendix B.2), so one labor-market condition is implied by the others.
We set $w_1=1$.

\paragraph{Existence and uniqueness.} We establish the results layer by layer: the
equilibrium of a single market, incomes given wages, and wages.

\begin{proposition}[One market]\label{prop:market}
For any positive delivered costs, a market has exactly one equilibrium.
\end{proposition}

\noindent The proof uses \eqref{eq:clearing}: each $\psi^{-1}(K_g/A)$ falls as $A$
rises, so the sum of shares is a strictly decreasing function of a single number and
equals one exactly once. Each group's profit is strictly concave in its own quantities,
so the first-order conditions identify the unique best response.

\begin{proposition}[Incomes]\label{prop:incomes}
Given wages, incomes are the unique solution of the linear system
\eqref{eq:income}, \eqref{eq:spending}, and $X_n\ge w_nL_n$ for every $n$.
\end{proposition}

\noindent Given wages, every share and markup is known, profits and tariff revenue are
linear in market sizes, and market sizes are linear in incomes. The system has a unique
solution because each dollar a country spends comes back to the world as profits and
tariff revenue by less than a dollar: the rest is paid to labor.

\begin{proposition}[Existence of wages]\label{prop:exist}
Suppose no set of countries is cut off from the rest: for every proper subset $Z$ of
countries, some plant located in $Z$ serves a market of a country outside $Z$. Then
there exist wages that clear every labor market.
\end{proposition}

\noindent Uniqueness of wages is the most delicate layer. Let $W_n(w)$ be the right
side of \eqref{eq:labor}, the wage bill of country $n$ at wages $w$, and
$M_{nj}\equiv\partial\ln W_n/\partial\ln w_j$ its elasticity with respect to country
$j$'s wage. Homogeneity makes every row of $M$ sum to one. Labor in different countries
is a gross substitute when $M_{nj}>0$ for $n\neq j$: a rise in any other country's wage
raises country $n$'s wage bill. Let $b^n_{i,n'}$ be the share of country $n$'s wage bill
paid by plant $i$ out of its costs in country $n'$, and
\begin{equation}
  \bar E^{(j)}_n\equiv\sum_{i,n'}b^n_{i,n'}\,\frac{\partial\ln E_{n'k(i)}}{\partial\ln w_j},
  \label{eq:demandterm}
\end{equation}
the average response of the markets in which country $n$'s plants sell.

\begin{theorem}[Wages]\label{thm:wages}
Let $D$ be a set of wage vectors that is convex in logarithms: the straight line in
logarithms between any two of its points stays in $D$. Suppose that at every $w\in D$
(i) no group holds more than half of any market and (ii) $\bar E^{(j)}_n>0$ for all
$n\neq j$. Then $M_{nj}\ge\bar E^{(j)}_n>0$ for $n\neq j$ throughout $D$, and $D$
contains at most one equilibrium wage vector up to scale.
\end{theorem}

\noindent The bound $M_{nj}\ge\bar E^{(j)}_n$ says that a rise in a foreign wage can
lower country $n$'s wage bill only by shrinking the markets in which its plants sell;
every reallocation of sales between plants and between groups favors country $n$.
Condition (i) is what signs the effect of variable markups on labor demand. With gross
substitutes, a second equilibrium is ruled out by the classical argument of
\citet{ArrowBlockHurwicz1959}. Condition (ii) is not a theorem of the model. A country
whose groups lose profits when a foreign wage rises spends less, and if the markets
where country $n$ sells shrink enough, its wage bill can fall. That income effect is the
ownership channel the model is built to study. Both conditions are inequalities in
objects the model computes, and we report them with the quantification.

\paragraph{Solution algorithm.} The layers above map into nested loops. Given wages,
each market is solved for $A$ from \eqref{eq:clearing}, and incomes come from one
linear solve. Wages are updated by the damped tâtonnement
$w_n\leftarrow w_n^{1-\lambda}(W_n(w)/L_n)^{\lambda}$, with $\lambda\in(0,1)$, which raises
a country's wage when its labor is in excess demand and whose fixed points are the
equilibrium wages (Appendix~\ref{app:algo}).

\subsection{Comparative Statics}
\label{sec:cs}

\paragraph{Incidence within a market.} Consider a market with given plants and
spending, and a change in delivered costs. Let $\hat c_g=\sum_{i\in T_g}(s_i/S_g)\,
d\ln c_i$ be the average cost change of group $g$, $\varepsilon_\mu(S)\equiv
S\mu'(S)/\mu(S)=S/(1-S)$ the elasticity of the markup, and
\begin{equation}
  \Lambda_g\equiv\frac{1}{1+(\sigma-1)\varepsilon_\mu(S_g)}=\frac{1-S_g}{1+(\sigma-2)S_g},
  \qquad
  \varpi_g\equiv\frac{S_g\Lambda_g}{\sum_h S_h\Lambda_h}.
  \label{eq:incidenceweights}
\end{equation}
$\Lambda_g$ falls from one toward zero as the share rises, and the weights $\varpi_g$
add up to one.

\begin{proposition}[Incidence]\label{prop:incidence}
Holding spending fixed,
\begin{equation}
  d\ln S_g=-(\sigma-1)\Lambda_g\bigl(\hat c_g-d\ln P\bigr),
  \qquad
  d\ln P=\sum_g\varpi_g\,\hat c_g .
  \label{eq:incidence}
\end{equation}
If only group $j$'s costs change, its sales-weighted average price rises by the
fraction $\rho_j=1-(\sigma-1)\varepsilon_\mu(S_j)\Lambda_j(1-\varpi_j)\in(0,1]$ of its
average cost increase $\hat c_j$, and every rival $h$ gains share,
$d\ln S_h=(\sigma-1)\Lambda_h\varpi_j\hat c_j$, and raises its price.
\end{proposition}

\noindent In the limit of small groups, which is monopolistic competition,
$\Lambda_g=1$, $\varpi_g=S_g$ and $\rho_j=1$. Large groups instead absorb part of a cost
increase in their markup, so pass-through falls with the group's share, and since
$\varpi_g<S_g$ for the larger groups, the price index moves by less than their share
would suggest. Rivals are not passive either: a cost increase for one group raises
their shares, markups and prices. Because tariffs depend on where goods are produced, a
tariff on goods from $\ell$ raises $\hat c_g$ only by the share of group $g$'s sales
that comes from its plants in $\ell$: a group able to serve the market from several
countries is partly shielded, while a stand-alone exporter from $\ell$ bears the tariff
in full.

\paragraph{Ownership and the level of profit income.} Consider a country $H$ with
ownership shares $\theta_g\equiv\theta_{gH}$, and its profit income from a market,
$\pi^H=\sum_g\theta_g\pi_g$.

\begin{proposition}[Ownership-weighted concentration]\label{prop:ownership}
\begin{equation}
  \frac{\pi^H}{E}=\frac{1}{\sigma}\Bigl[\bar\theta
  +\kappa\bigl(\bar\theta\cdot\mathrm{HHI}+\mathrm{Cov}_S(\theta,S)\bigr)\Bigr],
  \label{eq:decomp}
\end{equation}
where $\bar\theta=\sum_gS_g\theta_g$ is the sales-weighted ownership share,
$\mathrm{HHI}=\sum_gS_g^2$, and
$\mathrm{Cov}_S(\theta,S)=\sum_g\theta_gS_g^2-\bar\theta\cdot\mathrm{HHI}$ is the
covariance between ownership and share computed with sales weights. A calculation that
applies $\bar\theta$ to total profit in the market, $\bar\theta\sum_g\pi_g$, misses
exactly $(E\kappa/\sigma)\,\mathrm{Cov}_S(\theta,S)$.
\end{proposition}

\noindent The first term is what a model with constant markups delivers: ownership
times revenue over $\sigma$. The second is an ownership-weighted Herfindahl index. Its
component $\bar\theta\cdot\mathrm{HHI}$ can be built from a country's ownership share and
a concentration index; the covariance cannot, since it requires knowing which groups a
country owns and how large each one is. Because profits are convex in shares, two
countries with the same share of a market's sales earn different profits if one owns
the large groups and the other the small ones. With constant markups, as under
monopolistic competition, $\kappa=0$ and a country's profit income is its ownership
share times total profit, whatever the distribution of ownership across groups. Summed
over markets with spending weights, the same algebra splits the economy-wide correction
into a covariance across markets between ownership shares and concentration, and the
average of the covariances within markets.

\paragraph{Ownership and the response of profit income.} Welfare depends on how profit
income responds to policy, and the response is governed by a different statistic.

\begin{proposition}[Response of owned profits]\label{prop:response}
Holding spending fixed, a change in the delivered costs of group $j$ alone changes
country $H$'s profit income by
\begin{equation}
  \frac{d\pi^H}{E}=\frac{\sigma-1}{\sigma}\,\varpi_j\,\hat c_j
  \sum_{h\neq j}S_h\Lambda_h\Bigl[\theta_h\bigl(1+2\kappa S_h\bigr)
  -\theta_j\bigl(1+2\kappa S_j\bigr)\Bigr].
  \label{eq:response}
\end{equation}
\end{proposition}

\noindent A cost increase for group $j$ alone, for example a tariff on goods from a
country from which only $j$ serves the market, reallocates sales from $j$ to its rivals
in proportion to $S_h\Lambda_h$, and each unit of share is worth $\theta(1+2\kappa S)$ to
country $H$, up to the factor $E/\sigma$. The country gains when the rivals it owns are
worth more to it at the margin than the group whose costs rose. The weights are
$S_h\Lambda_h$ rather than $S_h^2$, so large groups count for less in the response than
in the level, and what matters is the ownership of the affected group relative to its
rivals, not the covariance that governs the level. The ownership-weighted Herfindahl
index measures a country's stake in a market but is not a sufficient statistic for how
that stake responds to a tariff; in general equilibrium the response also runs through
wages and market sizes.

\paragraph{Wages and prices.} A rise in country $j$'s wage raises the delivered cost of
every plant located in $j$ one for one. By Proposition~\ref{prop:incidence}, it raises
the price index of every market by $\sum_g\varpi_g\omega_{gj}$, where $\omega_{gj}$ is
the share of group $g$'s sales in the market that comes from its plants in $j$: a
proportion between zero and one, so markups dampen a wage change but cannot reverse or
amplify it. Who gains from a fall in $j$'s wage depends on ownership. The fall lowers
the costs of every group with plants in $j$, and part of the gain accrues as profit to
the owners of those groups, wherever they reside. A policy that lowers a foreign wage
can therefore raise the income of the country that owns the plants located there. This
is the general-equilibrium channel through which ownership matters for trade policy.

\subsection{Extension: Endogenous Entry}
\label{sec:entry}

The baseline takes as given which plants each group operates and which markets they
serve. This subsection lets groups choose them, so that which affiliates exist, and
where each of them sells, is an equilibrium outcome.

\paragraph{Fixed costs and plant sets.} Each group now has a set of potential plants
$\cP_g$. A plant that serves market $(n,k)$ pays a fixed cost
\begin{equation}
  F_{i,n}=f_k\,w_\ell\,f^d_{\ell n},
  \label{eq:fixed}
\end{equation}
in units of labor of the country $\ell$ where it is located, with $f_k>0$ and
$f^d_{\ell n}\ge1$. The cost is incurred once per market served, whatever the quantity
sold, and does not depend on who owns the plant. Because it is paid in labor, it appears
as wage income and leaves the model homogeneous of degree one in wages. Every fixed cost
is specific to a plant--market pair, so a group's problem separates across
destinations: in each market $(n,k(g))$ it chooses the subset $T_{g,n}\subseteq\cP_g$ of
its plants that serve $n$.\footnote{A plant-level setup cost would link a group's
choices across markets and turn its problem into a combinatorial one over locations and
markets jointly \citep{Tintelnot2017}.} This is a simplified version of the multi-plant
location problem of \citet{Yang2023}.

\paragraph{The entry problem.} A group's plant set matters for its profit only through
its strength. For a set of plants $T\subseteq\cP_g$, let $K(T)=\sum_{i\in T}
c_i^{1-\sigma}$ be its strength and $\Phi(T)=\sum_{i\in T}F_{i,n}$ its fixed cost.
Given $A$, serving the market with $T$ yields
\begin{equation}
  U_g(T;A)=\Omega\bigl(K(T);A\bigr)-\Phi(T),
  \qquad
  \Omega(K;A)\equiv\frac{E}{\sigma}\bigl(S+\kappa S^2\bigr)\Big|_{S=\psi^{-1}(K/A)} .
  \label{eq:entrypayoff}
\end{equation}
Because $\Omega$ is strictly increasing in $K$, a set can be optimal only if no other set
delivers at least as much strength at no greater fixed cost. These efficient sets form a
chain, ordered by strength and by cost together, so the group's choice amounts to how
far to climb a ladder that runs from small and cheap to large and expensive. When a
group's plants pay the same fixed cost to serve the market, each rung adds the plant
with the next-lowest delivered cost, and the group's choice is a cutoff on its own list
of plants, the within-group analogue of the cutoff in \citet{GaubertItskhoki2021}.

\begin{definition}[Entry equilibrium in a market]\label{def:entry}
Given delivered costs, fixed costs and spending, an entry equilibrium is a collection of
plant sets $\{T_g\}$ and an aggregate $A$ such that (i) each $T_g$ maximizes
$U_g(\cdot\,;A)$ over subsets of $\cP_g$, and (ii) $A$ solves \eqref{eq:clearing} with
$K_g=K(T_g)$.
\end{definition}

\noindent Groups evaluate entry at the prevailing degree of competition, as in a
free-entry condition, while taking into account their effect on $A$ when they set
quantities. This notion turns each group's entry problem into the choice of one number
given one number.

\begin{theorem}[Entry]\label{thm:entry}
Let $\sigma\ge2$.
\begin{enumerate}[label=(\alph*)]
\item $\Omega(K;A)$ is strictly concave in $K$ at every share.
\item $\Omega_{KA}<0$ whenever the group's share is at most one half.
\item Suppose that, for any two candidate equilibria, no group would hold more than
half of the market if it operated the plants it chooses in one of them while facing the
aggregate of the other, and that no group is indifferent between two rungs of its
ladder. Then a market has at most one entry equilibrium.
\end{enumerate}
\end{theorem}

\noindent Here $\Omega_{KA}$ is the cross partial derivative of $\Omega$. Part (a) says
that strength has diminishing returns for the group, and part (b) that it is worth less
in a more competitive market. Together they imply that each group's optimal strength
falls as $A$ rises, while $A$ rises with every group's strength, so two equilibria with
different aggregates cannot both exist. No order of entry is imposed, and groups still
internalize competition among their own plants. The hypothesis is the same one-half
condition as in Theorem~\ref{thm:wages}.

\paragraph{General equilibrium with entry.} Fixed costs are paid to labor where the
plant is located, so they are added to the right side of \eqref{eq:labor}, and profits
distributed through $\theta$ are net of them. The results of Section~\ref{sec:ge} hold
for given plant sets, with one change: existence also requires fixed costs small enough
that every country's income is at least a fixed fraction of its labor income. Two
limits are new. First, plant sets change discretely with wages, so uniqueness at each
layer does not make the joint equilibrium of wages and plant sets unique. Second,
entry is chosen in whole plants, so an exact equilibrium in pure strategies need not
exist. The solution algorithm adds an outer loop over plant sets
(Appendix~\ref{app:algo}).

\setlength{\bibsep}{1pt}
\begin{thebibliography}{99}

\bibitem[Antr\`as(2003)]{Antras2003}
Antr\`as, Pol. ``Firms, Contracts, and Trade Structure.'' \emph{Quarterly Journal of
Economics} 118(4): 1375--1418, 2003.

\bibitem[Antr\`as and Helpman(2004)]{AntrasHelpman2004}
Antr\`as, Pol, and Elhanan Helpman. ``Global Sourcing.'' \emph{Journal of Political
Economy} 112(3): 552--580, 2004.

\bibitem[Arrow, Block, and Hurwicz(1959)]{ArrowBlockHurwicz1959}
Arrow, Kenneth J., H. D. Block, and Leonid Hurwicz. ``On the Stability of the
Competitive Equilibrium, II.'' \emph{Econometrica} 27(1): 82--109, 1959.

\bibitem[Atkeson and Burstein(2008)]{AtkesonBurstein2008}
Atkeson, Andrew, and Ariel Burstein. ``Pricing-to-Market, Trade Costs, and
International Relative Prices.'' \emph{American Economic Review} 98(5): 1998--2031,
2008.

\bibitem[Gaubert and Itskhoki(2021)]{GaubertItskhoki2021}
Gaubert, Cecile, and Oleg Itskhoki. ``Granular Comparative Advantage.'' \emph{Journal
of Political Economy} 129(3): 871--939, 2021.

\bibitem[Head and Mayer(2019)]{HeadMayer2019}
Head, Keith, and Thierry Mayer. ``Brands in Motion: How Frictions Shape Multinational
Production.'' \emph{American Economic Review} 109(9): 3073--3124, 2019.

\bibitem[Mas-Colell, Whinston, and Green(1995)]{MWG1995}
Mas-Colell, Andreu, Michael D. Whinston, and Jerry R. Green. \emph{Microeconomic
Theory}. Oxford University Press, 1995.

\bibitem[Milgrom and Shannon(1994)]{MilgromShannon1994}
Milgrom, Paul, and Chris Shannon. ``Monotone Comparative Statics.''
\emph{Econometrica} 62(1): 157--180, 1994.

\bibitem[Ramondo and Rodr\'iguez-Clare(2013)]{RamondoRodriguezClare2013}
Ramondo, Natalia, and Andr\'es Rodr\'iguez-Clare. ``Trade, Multinational Production,
and the Gains from Openness.'' \emph{Journal of Political Economy} 121(2): 273--322,
2013.

\bibitem[Tintelnot(2017)]{Tintelnot2017}
Tintelnot, Felix. ``Global Production with Export Platforms.'' \emph{Quarterly Journal
of Economics} 132(1): 157--209, 2017.

\bibitem[Yang(2023)]{Yang2023}
Yang, Chenying. ``Location Choices of Multi-plant Oligopolists: Theory and Evidence
from the Cement Industry.'' Working paper, 2023.

\end{thebibliography}

\clearpage
\section*{Appendix}
\renewcommand{\theHsubsection}{app.\Alph{subsection}}
\renewcommand{\theHequation}{app.\arabic{equation}}
\setcounter{subsection}{0}
\renewcommand{\thesubsection}{\Alph{subsection}}

\subsection{The Market Game}
\label{app:market}

\paragraph{A.1 Proof of Proposition~\ref{prop:markup}.} Fix a market. Let
$\gamma\equiv(\sigma-1)/\sigma$ and $Q\equiv(\sum_iq_i^{\gamma})^{1/\gamma}$, the
quantity of the sectoral composite. Since $PQ=E$, inverse demand is
$p_i=E\,Q^{-\gamma}q_i^{\gamma-1}$, so $s_i=q_i^\gamma/Q^\gamma$ and group $g$'s revenue
is $E\,Q^{-\gamma}\sum_{j\in T_g}q_j^{\gamma}=E\,S_g$. Using
$\partial\ln Q/\partial q_i=q_i^{\gamma-1}/Q^{\gamma}=s_i/q_i$, the derivative of the
group's revenue with respect to $q_i$ is
\[
  E\,Q^{-\gamma}\Bigl[\gamma\,q_i^{\gamma-1}-\gamma\,\frac{q_i^{\gamma-1}}{Q^\gamma}
  \sum_{j\in T_g}q_j^\gamma\Bigr]
  =\gamma\,p_i\,(1-S_g).
\]
The second term in the bracket is the loss the group imposes on its own plants. Setting
marginal revenue equal to $c_i$ gives \eqref{eq:markup}. \qed

\paragraph{A.2 Proof of Proposition~\ref{prop:market}.} \emph{Best responses.} Given
rivals' quantities, group $g$'s profit is $E\,X_g/(X_g+B_g)-\sum_{i\in T_g}c_iq_i$,
where $X_g=\sum_{i\in T_g}q_i^\gamma$ and $B_g>0$ is the same index for rivals. The
function $X\mapsto X/(X+B)$ is strictly increasing and strictly concave, with
derivatives $B/(X+B)^2>0$ and $-2B/(X+B)^3<0$, and $X_g$ is strictly concave in
$(q_i)_{i\in T_g}$. A strictly increasing, strictly concave function of a strictly
concave function is strictly concave, and cost is linear, so profit is strictly concave
in the group's own quantities. Since $\partial X_g/\partial q_i\to\infty$ as $q_i\to0$,
the optimum is interior and is the unique solution of the first-order conditions.
\emph{The share map.} By \eqref{eq:sharemap}, $\psi(0)=0$, $\psi$ is continuous and
strictly increasing on $[0,1)$, and $\psi(S)\to\infty$ as $S\to1$, so it has a
continuous, strictly increasing inverse defined on $[0,\infty)$. \emph{The aggregate.}
By Proposition~\ref{prop:markup} and \eqref{eq:sharemap}, an equilibrium is an $A$
solving \eqref{eq:clearing}, and conversely any such $A$ generates prices
$p_i=\mu(S_g)c_i$ that satisfy every group's first-order conditions and are consistent
with $A$, since $\sum_ip_i^{1-\sigma}=\sum_g\mu(S_g)^{1-\sigma}K_g=A\sum_gS_g=A$. The
left side of \eqref{eq:clearing}, $\mathcal S(A)$, is continuous and strictly
decreasing. As $A\to\infty$ every share tends to zero, and as $A\to0$ every share tends
to one, so $\mathcal S$ tends to the number of groups, which is at least two. Hence
$\mathcal S(A)=1$ has exactly one solution. \qed

\subsection{General Equilibrium}
\label{app:ge}

\paragraph{B.1 Proof of Proposition~\ref{prop:incomes}.} Given wages, delivered costs
are known, and by Proposition~\ref{prop:market} so are all shares and markups. In
market $(n',k)$, group $g$'s profit is $E_{n'k}S_{g,n'}m(S_{g,n'})$ and tariff revenue is
$\tau_{n'k}E_{n'k}$, with $\tau_{n'k}=\sum_{i\in\cA_{n'k}}t_is_i/[\mu_{g(i)}(1+t_i)]$
and $t_i$ the tariff on plant $i$. Substituting \eqref{eq:spending} into
\eqref{eq:income},
\[
  X_n=w_nL_n+\sum_{n'}\Gamma_{nn'}X_{n'},
  \qquad
  \Gamma_{nn'}=\sum_k\beta_k\Bigl[\sum_{g:\,k(g)=k}\theta_{gn}S_{g,n'}m(S_{g,n'})
  +\mathbf 1\{n'=n\}\,\tau_{nk}\Bigr]\ge0 .
\]
Using $\sum_n\theta_{gn}=1$, the column of $\Gamma$ for $n'$ sums to
$\sum_k\beta_k[\sum_gS_{g,n'}m(S_{g,n'})+\tau_{n'k}]$. Per dollar of spending in a
market, revenue splits into profit, tariff revenue and labor cost,
\begin{equation}
  \sum_gS_gm(S_g)+\tau+\sum_i\frac{s_i}{\mu_{g(i)}(1+t_i)}=1,
  \tag{B.1}\label{eq:split}
\end{equation}
and the last term is positive. So every column of $\Gamma$ sums to less than one, the
largest eigenvalue of $\Gamma$ in absolute value is below one, and
$X=(I-\Gamma)^{-1}wL=\sum_{s\ge0}\Gamma^swL\ge wL$ is the unique solution. \qed

\paragraph{B.2 Proof of Proposition~\ref{prop:exist}.} Normalize wages to the unit
simplex and let $z_n(w)=W_n(w)/w_n-L_n$ be excess demand for labor. By
Propositions~\ref{prop:market}--\ref{prop:incomes}, $z$ is continuous on the interior of
the simplex and homogeneous of degree zero. It satisfies Walras' law
$\sum_nw_nz_n(w)=0$. Adding \eqref{eq:income} across countries, world income is
$\sum_nw_nL_n$ plus world profits plus world tariff revenue. By \eqref{eq:spending},
world spending equals world income, and by \eqref{eq:split} it also equals
$\sum_nW_n$ plus world profits plus world tariff revenue. Hence
$\sum_nW_n=\sum_nw_nL_n$. Excess demand is bounded below by $-L$. It remains to check
the boundary condition of \citet[Proposition 17.C.1]{MWG1995}: if $w^m\to\bar w$, where
$Z=\{n:\bar w_n=0\}$ is non-empty, then $\max_nz_n(w^m)\to\infty$.

A plant's delivered cost is proportional to the wage where it is located. So the costs
of plants located in $Z$ go to zero, while all other costs stay in a compact subset of
$(0,\infty)$. By hypothesis, some plant located in $Z$ serves a market $(n',k)$ with
$n'\notin Z$. By Proposition~\ref{prop:incomes}, $X_{n'}\ge w_{n'}L_{n'}$, so
$E_{n'k}=\beta_kX_{n'}$ stays bounded away from zero. For $m$ large, the lowest-cost
plant $i$ in that market is located in some $\ell\in Z$. Let $g$ be its group and $J$
the number of plants in the market. Its cost is $V_i=E_{n'k}s_i/[\mu_g(1+t_i)]$, all of
it paid to labor in $\ell$, so it suffices to show that $V_i/w_\ell\to\infty$.

If $S_g\le\frac12$, then $\mu_g\le2\sigma/(\sigma-1)$. Every price is at least $c_i$,
since markups exceed one, so $P^{1-\sigma}\le Jc_i^{1-\sigma}$ and
$s_i=(\mu_gc_i/P)^{1-\sigma}\ge\mu_g^{1-\sigma}/J$. Then $V_i$ is bounded away from zero
while $w_\ell\to0$.

If $S_g>\frac12$, every rival $h$ has $S_h\le1-S_g<\frac12$, so
$(1-S_h)^{1-\sigma}\le2^{\sigma-1}$ and $1/S_g<2$. By \eqref{eq:sharemap},
\[
  \frac{K_h}{K_g}=\frac{\psi(S_h)}{\psi(S_g)}
  =\frac{S_h(1-S_h)^{1-\sigma}}{S_g(1-S_g)^{1-\sigma}}
  \le2^{\sigma}(1-S_g)^{\sigma},
\]
so $1-S_g\ge\frac12(K_h/K_g)^{1/\sigma}$. Rival strength $K_h$ is bounded below, because
all costs are bounded above. Plant $i$ is the cheapest in its group, so
$K_g\le Jc_i^{1-\sigma}$ and $s_i\ge S_g/J>1/(2J)$. With $1/\mu_g=(\sigma-1)(1-S_g)/\sigma$,
this gives $V_i\ge c\,c_i^{(\sigma-1)/\sigma}$ for some $c>0$, and since $c_i$ is
proportional to $w_\ell$, $V_i/w_\ell$ grows like $w_\ell^{-1/\sigma}$.

In either case $z_\ell\to\infty$, and the result follows from
\citet[Proposition 17.C.1]{MWG1995}. \qed

\paragraph{B.3 A decomposition of labor payments.} Raise wages, holding tariffs and
trade costs fixed, and let hats denote log changes. A plant's delivered cost changes by
$\hat c_i=\hat w_{\ell(i)}$. Its cost in a market is $V_i=r_i/[\mu_g(1+t_i)]$, so
$\hat V_i=\hat E+\hat s_i-\hat\mu_g$ with $\hat s_i=(1-\sigma)(\hat c_i+\hat\mu_g-\hat P)$
and $\hat\mu_g=\varepsilon_\mu\hat S_g$. By Proposition~\ref{prop:incidence},
$\hat S_g=-(\sigma-1)\Lambda_g(\hat c_g-\hat P)$, and using
$1=\Lambda_g[1+(\sigma-1)\varepsilon_\mu]$,
\begin{equation}
  \hat V_i=\hat E-(\sigma-1)(\hat c_i-\hat c_g)-\chi_g\,(\hat c_g-\hat P),
  \qquad
  \chi_g\equiv(1-\varepsilon_\mu)(\sigma-1)\Lambda_g
  =\frac{(\sigma-1)(1-2S_g)}{1+(\sigma-2)S_g} .
  \tag{B.3}\label{eq:vdecomp}
\end{equation}
The first term is the growth of the market, the second the reallocation of sales within
the group, the third the reallocation between the group and its rivals. The coefficient
$\chi_g$ lies in $[0,\sigma-1]$ exactly when $S_g\le\frac12$.

\paragraph{B.4 Proof of Theorem~\ref{thm:wages}.} \emph{The bound.} Raise $w_j$ alone,
$j\neq n$. Every plant paying country $n$ is located in $n$, so $\hat c_i=0$ for those
plants. Every other plant has $\hat c_i\in\{0,1\}$, so $\hat c_g\in[0,1]$ and, by
Proposition~\ref{prop:incidence}, $\hat P=\sum_g\varpi_g\hat c_g\in[0,1]$. With
$\hat c_i=0$ the last two terms of \eqref{eq:vdecomp} equal
$(\sigma-1)\hat c_g-\chi_g(\hat c_g-\hat P)$. If $\hat c_g\ge\hat P$, condition (i)
gives $\chi_g\le\sigma-1$ and the expression is at least $(\sigma-1)\hat P\ge0$;
otherwise both parts are non-negative, since $\chi_g\ge0$. So $\hat V_i\ge\hat E$ for
every plant paying country $n$, and weighting by $b^n_{i,n'}$, which sum to one, gives
$M_{nj}\ge\bar E^{(j)}_n$. \emph{Uniqueness.} The argument is classical
\citep{ArrowBlockHurwicz1959}; see also \citet[Section 17.F]{MWG1995}. Since $M$ is
homogeneous of degree zero in wages, conditions (i) and (ii) also hold at every
rescaling of a point of $D$, and the set of these rescalings is still convex in
logarithms. Let $w,w^*\in D$ be non-proportional equilibria. Scale $w^*$ so that
$w^*\le w$ with equality at some $n_0$, and move from $w^*$ to $w$ along the segment in
logarithms. Every wage other than $w_{n_0}$ weakly rises and at least one strictly
rises, so $W_{n_0}$ strictly rises while $w_{n_0}$ does not. Then
$z_{n_0}(w)>z_{n_0}(w^*)=0$, a contradiction. \qed

\subsection{Comparative Statics}
\label{app:cs}

\paragraph{C.1 Proof of Proposition~\ref{prop:incidence}.} By \eqref{eq:strength},
$d\ln K_g=(1-\sigma)\hat c_g$, and $d\ln A=(1-\sigma)\,d\ln P$. Differentiating
$\psi(S_g)=K_g/A$, with $d\ln\psi/d\ln S=1+(\sigma-1)\varepsilon_\mu=1/\Lambda_g$, gives
$d\ln S_g=-(\sigma-1)\Lambda_g(\hat c_g-d\ln P)$. Shares sum to one, so
$\sum_gS_g\,d\ln S_g=0$, which gives $d\ln P=\sum_g\varpi_g\hat c_g$. If only $j$'s
costs change, $d\ln P=\varpi_j\hat c_j$, so $d\ln S_j=-(\sigma-1)\Lambda_j(1-\varpi_j)\hat
c_j$ and, for $h\neq j$, $d\ln S_h=(\sigma-1)\Lambda_h\varpi_j\hat c_j$. Prices are
$p_i=\mu(S_g)c_i$, so $d\ln p_i=d\ln c_i+\varepsilon_\mu(S_g)\,d\ln S_g$. Averaging over
group $j$'s plants with weights $s_i/S_j$ gives $\rho_j\hat c_j$, and every rival's
price rises. Since $(\sigma-1)\varepsilon_\mu\Lambda_j<1$ and $1-\varpi_j\in[0,1)$,
$\rho_j\in(0,1]$. \qed

\paragraph{C.2 Proof of Proposition~\ref{prop:ownership}.} By \eqref{eq:profit},
$\pi^H/E=\sigma^{-1}\sum_g\theta_gS_g(1+\kappa S_g)=\sigma^{-1}[\bar\theta+\kappa\sum_g
\theta_gS_g^2]$, and $\sum_g\theta_gS_g^2=\bar\theta\cdot\mathrm{HHI}+\mathrm{Cov}_S
(\theta,S)$ by definition. Total profit is the case $\theta_g\equiv1$,
$\sum_g\pi_g/E=(1+\kappa\,\mathrm{HHI})/\sigma$, using $\sum_gS_g=1$. Subtracting
$\bar\theta\sum_g\pi_g$ from $\pi^H$ leaves $(E\kappa/\sigma)\mathrm{Cov}_S(\theta,S)$.
\qed

\paragraph{C.3 Proof of Proposition~\ref{prop:response}.} By \eqref{eq:profit},
$d\pi_g=(E/\sigma)(1+2\kappa S_g)\,dS_g$. Let $Z=\sum_hS_h\Lambda_h$, so
$\varpi_j=S_j\Lambda_j/Z$. From C.1,
$dS_j=-(\sigma-1)S_j\Lambda_j(1-\varpi_j)\hat c_j=-(\sigma-1)\varpi_j\hat c_j\sum_{h\neq
j}S_h\Lambda_h$ and $dS_h=(\sigma-1)\varpi_jS_h\Lambda_h\hat c_j$ for $h\neq j$.
Substituting into $d\pi^H=\sum_g\theta_gd\pi_g$ and collecting terms gives
\eqref{eq:response}. \qed

\subsection{Entry}
\label{app:entry}

\paragraph{D.1 Two elasticities.} Let $G(S)\equiv\pi'(S)/\psi'(S)$, where
$\pi(S)=E(S+\kappa S^2)/\sigma$, and define $\varepsilon_G=SG'/G$ and
$\varepsilon_\psi=S\psi'/\psi$. Differentiating $\psi(S)=K/A$ gives
$\partial S/\partial K=1/(A\psi')$ and $\partial S/\partial A=-\psi/(A\psi')$. Hence
$\Omega_K=G(S)/A>0$, $\Omega_{KK}=G'(S)/(A^2\psi'(S))$, and
\[
  \Omega_{KA}=-\frac{1}{A^2}\Bigl[G'(S)\frac{\psi(S)}{\psi'(S)}+G(S)\Bigr]
  =-\frac{G(S)}{A^2\varepsilon_\psi(S)}\bigl[\varepsilon_G(S)+\varepsilon_\psi(S)\bigr].
\]
So $\Omega$ is strictly concave in $K$ if and only if $\varepsilon_G<0$, and
$\Omega_{KA}<0$ if and only if $\varepsilon_G+\varepsilon_\psi>0$.

\paragraph{D.2 Proof of Theorem~\ref{thm:entry}(a)--(b).} Set $E=1$,
$a\equiv2(\sigma-1)$ and $b\equiv\sigma-2\ge0$. Then $\pi'(S)=(1+aS)/\sigma$ and
$\psi'(S)=C(1-S)^{-\sigma}(1+bS)$ with $C=[\sigma/(\sigma-1)]^{\sigma-1}$, so
\[
  \varepsilon_G(S)=S\,f(S),\qquad
  f(S)=\frac{a}{1+aS}-\frac{b}{1+bS}-\frac{\sigma}{1-S}.
\]
At $S=0$, $f(0)=a-b-\sigma=0$. Since $a>b\ge0$, $a/(1+aS)>b/(1+bS)\ge0$, so
$f'(S)=-a^2/(1+aS)^2+b^2/(1+bS)^2-\sigma/(1-S)^2<0$. Hence $f<0$ and $\varepsilon_G<0$ on
$(0,1)$, which is (a). Next, $\varepsilon_\psi=(1+bS)/(1-S)$, and using
$1+bS-\sigma S=1-2S$,
\[
  \varepsilon_G+\varepsilon_\psi=\frac{aS}{1+aS}-\frac{bS}{1+bS}+\frac{1-2S}{1-S}.
\]
The map $x\mapsto x/(1+x)$ is increasing and $a>b$, so the first two terms sum to a
positive number; the third is non-negative when $S\le\frac12$. This is (b). At
$S=\frac12$ the expression equals $(\sigma-1)/\sigma-(\sigma-2)/\sigma=1/\sigma$. \qed

\paragraph{D.3 Proof of Theorem~\ref{thm:entry}(c).} \emph{Step 1: the ladder.} Since
$\Omega_K>0$, a set $T$ is optimal only if no $T'$ has $K(T')\ge K(T)$ and
$\Phi(T')\le\Phi(T)$ with one inequality strict. The remaining efficient sets can be
ordered so that $K$ and $\Phi$ rise together. \emph{Step 2: single crossing.} Take
$A<A'$ and two rungs with $K_1<K_2$, and let
$\Delta(x)=\int_{K_1}^{K_2}\Omega_K(K;x)\,dK-[\Phi_2-\Phi_1]$ be the net gain from the
larger rung at aggregate $x$. The share $\psi^{-1}(K/x)$ rises with $K$ and falls with
$x$, so on $[K_1,K_2]\times[A,A']$ it is largest at $(K_2,A)$. If
$\psi^{-1}(K_2/A)\le\frac12$, part (b) gives $\Omega_{KA}<0$ throughout, and $\Delta$
is strictly decreasing on $[A,A']$. \emph{Step 3: the aggregate.} Given strengths, the
aggregate $A(\{K_g\})$ solving \eqref{eq:clearing} is strictly increasing in each $K_g$:
raising $K_j$ at a given $A$ raises the left side, which is strictly decreasing in $A$.
\emph{Step 4.} Suppose $(A,\{K_g\})$ and $(A',\{K'_g\})$ are equilibria with $A<A'$,
and suppose $K'_g>K_g$ for some $g$. Apply Step 2 with $K_1=K_g$ and $K_2=K'_g$: by the
hypothesis, group $g$ would hold at most half of the market with the plants it chooses
at $A'$ while facing $A$, so $\psi^{-1}(K'_g/A)\le\frac12$ and $\Delta$ is strictly
decreasing on $[A,A']$. Optimality at $A'$ gives $\Delta(A')\ge0$, hence $\Delta(A)>0$,
so at $A$ the group strictly prefers $K'_g$ to $K_g$, contradicting optimality at $A$.
Hence $K'_g\le K_g$ for every $g$ \citep[the single-crossing argument of][]{MilgromShannon1994},
and by Step 3, $A'=A(\{K'_g\})\le A(\{K_g\})=A$, a contradiction. Equilibria with the
same aggregate have the same best responses, and with no ties the same plant sets.
\qed

\paragraph{D.4 The notion of entry equilibrium.} Definition~\ref{def:entry} evaluates
every alternative set at the prevailing aggregate $A$. In a two-stage game in which the
market is re-solved after a deviation, a group that adds a plant anticipates that $A$
rises, which lowers the value of the deviation; the two notions coincide for groups
with negligible shares and can differ for large ones. Theorem~\ref{thm:entry} applies to
the notion of Definition~\ref{def:entry}. Existence in pure strategies is not guaranteed
under either notion: a group's choice is discrete, so its optimal strength jumps with
$A$, and \eqref{eq:clearing} may have no solution with each group on a rung.

\subsection{Solution Algorithm}
\label{app:algo}

\paragraph{E.1 Markets.} Given an aggregate $A$, each group's share solves the scalar
equation \eqref{eq:sharemap}, and the market solves \eqref{eq:clearing} by safeguarded
Newton steps on $\ln A$, which never leave a bracket that contains the unique root.

\paragraph{E.2 Incomes and wages.} Given wages, incomes come from the linear solve of
B.1. Wages are updated by the damped tâtonnement of Section~\ref{sec:ge}, followed by
Newton steps once close to the solution. When labor demand is elastic, $M_{nn}<0$ and an
undamped update ($\lambda=1$) can overshoot; choosing $\lambda$ with
$\lambda(1-M_{nn})<1$ keeps the update monotone.

\paragraph{E.3 Plant sets (extension).} With entry, each group's best response at a
given $A$ is found by scanning its ladder, and $A$ by bisection. An outer loop updates
plant sets: at the equilibrium wages for the current sets, the groups with the largest
gains from changing their sets do so, and the procedure repeats until no group wishes to
move. Moving only the groups with the largest gains changes the path of the algorithm,
not its fixed points, and does not amount to an order of entry. When the loop does not
reach a configuration in which no group wishes to move, it returns the configuration
with the fewest such groups and reports how many there are. Because the joint
equilibrium is not known to be unique, we solve from several starting points.

\clearpage
\section*{Notation}
\label{app:notation}

\noindent Symbols are grouped by topic, and the last column gives the equation or
section where each is defined. A subscript
$n$ or $n'$ denotes a destination country, $\ell$ the country where a plant is located,
$k$ a sector, $g$, $h$ or $j$ a group, and $i$ a plant. Within a single market the
market's indices are dropped.

\small
\begin{longtable}{@{}p{0.17\linewidth}p{0.62\linewidth}p{0.15\linewidth}@{}}
\toprule
Symbol & Meaning & Where \\
\midrule
\endhead
\bottomrule
\endfoot
\multicolumn{3}{@{}l}{\emph{Agents and sets}}\\
$\cN$, $N$ & set of countries; number of countries & Agents \\
$\cK$ & set of sectors & Agents \\
$\cG$ & set of groups (a group is a global ultimate parent and all its plants) & Agents \\
$h(g)$, $k(g)$ & headquarters country and sector of group $g$ & Agents \\
$\cP_g$ & plants owned by group $g$ (potential plants in the extension) & Agents \\
$g(i)$, $\ell(i)$ & group and location of plant $i$ & Agents \\
$\cA_{nk}$ & plants that serve market $(n,k)$ & \eqref{eq:utility} \\
$T_g$, $T_{g,n}$ & plants of group $g$ that serve a given market & Section~\ref{sec:market} \\[3pt]
\multicolumn{3}{@{}l}{\emph{Households and demand}}\\
$L_n$ & labor supplied by country $n$ & Agents \\
$U_n$, $C_{nk}$ & utility; consumption of the sector-$k$ bundle & \eqref{eq:utility} \\
$\beta_k$ & share of income spent on sector $k$ & \eqref{eq:utility} \\
$\sigma_k$ & elasticity of substitution between varieties of sector $k$ & \eqref{eq:utility} \\
$q_{i,n}$, $p_{i,n}$ & quantity and price of plant $i$'s variety in country $n$ & \eqref{eq:demand} \\
$P_{nk}$ & price index of sector $k$ in country $n$ & \eqref{eq:demand} \\
$P_n$ & consumer price index of country $n$ & \eqref{eq:demand} \\
$X_n$ & income of country $n$ & \eqref{eq:income} \\
$E_{nk}$ & spending on market $(n,k)$, the size of the market & \eqref{eq:spending} \\[3pt]
\multicolumn{3}{@{}l}{\emph{Technology, trade and ownership}}\\
$w_\ell$ & wage in country $\ell$ & \eqref{eq:cost} \\
$\varphi_i$ & productivity of plant $i$ & \eqref{eq:cost} \\
$a_i$ & unit cost of plant $i$, $w_\ell/\varphi_i$ & \eqref{eq:cost} \\
$d_{\ell n}$ & iceberg trade cost from $\ell$ to $n$ ($\ge1$) & \eqref{eq:delivered} \\
$t_{\ell nk}$ & tariff of country $n$ on sector-$k$ goods produced in $\ell$ & \eqref{eq:delivered} \\
$c_{i,n}$ & delivered marginal cost of plant $i$ in market $n$, including the tariff & \eqref{eq:delivered} \\
$\theta_{gn}$ & share of group $g$ owned by country $n$ & \eqref{eq:income} \\
$\pi_g$, $\pi_{g,n}$ & profit of group $g$, in total and in market $n$ & \eqref{eq:income} \\
$\TR_n$ & tariff revenue of country $n$ & \eqref{eq:income} \\[3pt]
\multicolumn{3}{@{}l}{\emph{Market structure (one market)}}\\
$A$ & market aggregate, $P^{1-\sigma}$; high when the market is competitive & \eqref{eq:shares} \\
$s_i$ & revenue share of plant $i$ in the market & \eqref{eq:shares} \\
$S_g$ & revenue share of group $g$ in the market & \eqref{eq:shares} \\
$\mu(S)$ & markup of a group with share $S$ & \eqref{eq:markup} \\
$K_g$ & strength of group $g$, $\sum_{i\in T_g}c_i^{1-\sigma}$ & \eqref{eq:strength} \\
$\psi(S)$ & share map: a group with strength $K$ has the share $S$ with $\psi(S)=K/A$ & \eqref{eq:sharemap} \\
$m(S)$ & Lerner index, $1-1/\mu(S)$ & \eqref{eq:profit} \\
$\kappa$ & curvature of profit in the share, $\sigma-1$ & \eqref{eq:profit} \\[3pt]
\multicolumn{3}{@{}l}{\emph{General equilibrium}}\\
$V_{i,n'}$ & cost of plant $i$ in serving country $n'$, all paid to labor & \eqref{eq:labor} \\
$r_{i,n'}$ & revenue of plant $i$ in country $n'$ & \eqref{eq:labor} \\
$W_n$ & wage bill of country $n$ (labor demanded, in value) & Section~\ref{sec:ge} \\
$M_{nj}$ & elasticity of $W_n$ with respect to the wage of country $j$ & Section~\ref{sec:ge} \\
$b^n_{i,n'}$ & share of $W_n$ paid by plant $i$ out of its costs in $n'$ & \eqref{eq:demandterm} \\
$\bar E^{(j)}_n$ & average response of the markets where country $n$'s plants sell to $w_j$ & \eqref{eq:demandterm} \\
$D$ & set of wage vectors on which uniqueness is checked & Theorem~\ref{thm:wages} \\
$\lambda$ & damping of the wage update & Solution algorithm \\
$z_n$ & excess demand for labor, $W_n/w_n-L_n$ & Appendix B.2 \\
$\Gamma$, $\tau_{nk}$ & income matrix; tariff revenue per dollar spent in $(n,k)$ & Appendix B.1 \\[3pt]
\multicolumn{3}{@{}l}{\emph{Comparative statics}}\\
$\hat x$ & log change of $x$ & Section~\ref{sec:cs} \\
$\hat c_g$ & sales-weighted average change in group $g$'s delivered costs & \eqref{eq:incidenceweights} \\
$\varepsilon_\mu(S)$ & elasticity of the markup with respect to the share, $S/(1-S)$ & \eqref{eq:incidenceweights} \\
$\Lambda_g$ & how strongly group $g$'s share responds to its relative cost & \eqref{eq:incidenceweights} \\
$\varpi_g$ & weight of group $g$ in the response of the price index & \eqref{eq:incidenceweights} \\
$\rho_j$ & pass-through of a cost increase to group $j$'s prices & Proposition~\ref{prop:incidence} \\
$\pi^H$, $\theta_g$ & profit income of country $H$ from a market; its ownership shares & \eqref{eq:decomp} \\
$\bar\theta$ & sales-weighted ownership share of country $H$ in the market & \eqref{eq:decomp} \\
HHI & Herfindahl index of group shares, $\sum_gS_g^2$ & \eqref{eq:decomp} \\
$\mathrm{Cov}_S(\theta,S)$ & sales-weighted covariance between ownership and share & \eqref{eq:decomp} \\
$\omega_{gj}$ & share of group $g$'s sales in a market that comes from its plants in $j$ & Wages and prices \\
$\chi_g$ & coefficient of the between-group reallocation of labor payments & Appendix B.3 \\[3pt]
\multicolumn{3}{@{}l}{\emph{Extension: entry}}\\
$F_{i,n}$ & fixed cost for plant $i$ to serve market $n$ & \eqref{eq:fixed} \\
$f_k$, $f^d_{\ell n}$ & sector and distance components of the fixed cost & \eqref{eq:fixed} \\
$K(T)$, $\Phi(T)$ & strength and fixed cost of a set of plants $T$ & \eqref{eq:entrypayoff} \\
$\Omega(K;A)$ & profit of a group with strength $K$ at aggregate $A$ & \eqref{eq:entrypayoff} \\
$U_g(T;A)$ & net payoff of group $g$ from serving the market with $T$ & \eqref{eq:entrypayoff} \\
\end{longtable}

\end{document}
